%==============================================================================
\section*{target và ký hiệu}
\addcontentsline{toc}{section}{target và ký hiệu}
%==============================================================================

\begin{phuongphap}[title={target học tập}]
Sau chương này, người học có thể:
\begin{enumerate}
  \item Giải thích \textbf{bản chất} của Gradient Descent: đi theo hướng \textbf{ngược gradient} để giảm hàm target (thường là hàm mất mát).
  \item Viết và chạy GD cho hàm một biến và nhiều biến; kiểm tra gradient bằng \textbf{numerical gradient}.
  \item Nhận biết ảnh hưởng của \textbf{điểm khởi tạo}, \textbf{learning rate} và hình dạng đường đồng mức tới hội tụ.
  \item Phân biệt \textbf{Batch GD}, \textbf{SGD}, \textbf{mini-batch}; mô tả vai trò của \textbf{Momentum}, \textbf{NAG} và ý tưởng \textbf{Newton}.
  \item Chọn điều kiện dừng phù hợp trong thực nghiệm.
\end{enumerate}
\end{phuongphap}

\begin{luuy}[title={Kiến thức tiên quyết}]
Đạo hàm hàm một biến và nhiều biến; ma trận và nhân ma trận; Linear Regression (hàm mất mát dạng bình phương tối thiểu). Phần code dùng \texttt{numpy} và \texttt{matplotlib}.
\end{luuy}

\begin{phuongphap}[title={Ký hiệu thống nhất}]
\begin{itemize}
  \item $f$, $\mathcal{L}$, $J$: hàm target / hàm mất mát.
  \item $\boldsymbol{\theta}$ hoặc $\mathbf{w}$: vector tham số cần tối ưu.
  \item $\eta$: learning rate (bước học); $\gamma$: hệ số momentum.
  \item $\nabla f$, $\nabla_{\boldsymbol{\theta}}\mathcal{L}$: gradient (vector chỉ hướng \textbf{tăng nhanh nhất} của hàm).
  \item $x^*$, $\boldsymbol{\theta}^*$: điểm cực tiểu (local hoặc global tùy bài toán).
  \item GD: Gradient Descent; SGD: Stochastic Gradient Descent.
\end{itemize}
\end{phuongphap}

\begin{dongco}[title={Bản chất Gradient Descent}]
\textbf{Gradient Descent} không phải công thức đóng cho nghiệm tối ưu, mà là quy trình \textbf{lặp}:
\begin{enumerate}
  \item Tại tham số hiện tại, tính \textbf{gradient} của hàm mất mát.
  \item Cập nhật tham số một bước nhỏ theo hướng \textbf{trừ gradient} (vì gradient chỉ hướng \textbf{đi lên dốc nhanh nhất} của hàm).
  \item Dừng khi gradient đủ nhỏ hoặc đạt giới hạn vòng lặp.
\end{enumerate}
Trong Machine Learning, ta thường \textbf{không giải được} $\nabla \mathcal{L} = 0$ bằng đại số (hàm phi tuyến, số chiều lớn, dữ liệu lớn). GD thay thế bằng đi từng bước nhỏ --- giống ``lăn xuống thung lũng'' trên bề mặt mất mát.
\end{dongco}

\newpage
